% !TeX root = ../proyecto.tex

\section{Requerimientos No Funcionales}

Se detallan los requerimientos de calidad, seguridad, desempeño y tolerancia a fallos del
sistema inteligente.

\begin{NFRequerimientos}
    % \NFRitem{ID}{Atributo}{Necesidad}{Estrategia}{Referencia}

    \NFRitem{RNF-01}{Optimización de Memoria (OOM)}
            {El modelo BETO no debe provocar cierres abruptos del proceso por saturación de RAM
            durante la inferencia en lote.}
            {Implementar liberación explícita del caché de PyTorch y reducción dinámica del
            tamaño del lote (\textit{batch-size}) a la unidad (\texttt{batch\_size = 1}) ante
            excepciones \texttt{MemoryError} o \texttt{RuntimeError}. La técnica de
            \textit{Gradient Accumulation} permite simular lotes grandes usando RAM mínima
            durante el Fine-Tuning.}
            {CU-01}

    \NFRitem{RNF-02}{Evasión de Contramedidas Web (WAF)}
            {El recolector debe poder acceder a periódicos digitales protegidos por
            Web Application Firewalls empresariales (Cloudflare, Akamai) sin ser bloqueado
            sistemáticamente.}
            {Uso del módulo \texttt{StealthSession} v7.0 con emulación de fingerprint TLS via
            \texttt{cloudscraper}, delays con distribución log-normal ($\mu=2.5$, $\sigma=0.7$)
            y patrón \textit{Circuit Breaker} por dominio que evita el bloqueo permanente de IP
            tras 3 fallos consecutivos (cooldown de 300~s con extensión a 600~s).}
            {CU-01, CU-02}

    \NFRitem{RNF-03}{Resiliencia en APIs Externas}
            {El módulo \texttt{DeepInvestigator} no debe terminar abruptamente si DuckDuckGo
            bloquea las peticiones de scraping de búsqueda.}
            {Implementar fallback hacia la herramienta nativa \texttt{GoogleSearch} de Gemini,
            ejecutando la búsqueda directamente dentro del modelo sin scraping externo. Si
            ambos mecanismos fallan, el estado de la investigación pasa a \texttt{error}
            con notificación al usuario.}
            {CU-03}

    \NFRitem{RNF-04}{Integridad de Base de Datos (ACID)}
            {Los fallos de conexión o errores de validación durante inserciones masivas
            no deben dejar la base de datos en estados corruptos o con registros parciales.}
            {Uso estricto de transacciones atómicas de SQLAlchemy con \textit{Rollback}
            automático ante excepciones del tipo \texttt{OperationalError} o
            \texttt{IntegrityError}. Las restricciones \texttt{FOREIGN KEY} con
            \texttt{CASCADE DELETE} garantizan la coherencia referencial sin lógica
            adicional en la capa de aplicación.}
            {CU-01, CU-05}

    \NFRitem{RNF-05}{Procesamiento Asíncrono No Bloqueante}
            {Las operaciones de larga duración (pipeline de análisis y consultas a Gemini)
            no deben bloquear la interfaz web del Dashboard ni la atención de otras
            peticiones HTTP.}
            {Uso de hilos en segundo plano (\textit{daemon threads}) gestionados con
            \texttt{threading.Thread} y un mapa de estado en memoria
            (\texttt{\_analysis\_status}, \texttt{\_active\_investigations}) consultado
            por el frontend mediante \textit{polling} periódico a los endpoints
            \texttt{/api/analyze/status} y \texttt{/api/investigate/status/<id>}.}
            {CU-01, CU-03}

    \NFRitem{RNF-06}{Autenticación y Control de Acceso}
            {Todos los endpoints del Dashboard y la API REST deben requerir una sesión
            autenticada. El acceso sin autenticación debe redirigir al formulario de
            inicio de sesión.}
            {Implementación de Flask-Login con el decorador \texttt{@login\_required}
            en todas las rutas excepto \texttt{/health}, \texttt{/login} y
            \texttt{/register}. Las contraseñas se almacenan como hashes
            \texttt{bcrypt} y nunca en texto plano. La protección contra redirecciones
            abiertas (\textit{Open Redirect}) se valida en el helper
            \texttt{\_is\_safe\_url()}.}
            {CU-06}

    \NFRitem{RNF-07}{Compatibilidad Dual PostgreSQL / CSV}
            {El sistema debe mantener la capacidad de operar sobre datos en formato
            CSV cuando la base de datos PostgreSQL no esté disponible, garantizando
            continuidad operativa básica del Dashboard.}
            {Implementación del patrón \textit{Check-then-Fallback} en todos los
            endpoints de datos: \texttt{\_check\_postgres()} verifica si existen
            registros en PostgreSQL; en caso negativo, \texttt{\_load\_data()}
            carga el CSV del directorio \texttt{data/}. El campo
            \texttt{data\_source} en cada respuesta JSON indica el motor activo.}
            {CU-01, CU-07}

    \NFRitem{RNF-08}{Trazabilidad por Lote de Análisis (Batch ID)}
            {El sistema debe permitir consultar el historial de ejecuciones
            independientes del pipeline y filtrar los datos del Dashboard por
            lote específico para análisis comparativo.}
            {Cada ejecución del pipeline asigna un \texttt{batch\_id} único a
            todos los registros insertados en ese ciclo. Los endpoints
            \texttt{/api/stats}, \texttt{/api/noticias} y \texttt{/api/batches}
            aceptan el parámetro \texttt{batch\_id} (\texttt{latest},
            \texttt{all} o un identificador específico).}
            {CU-01}

\end{NFRequerimientos}
